\documentclass{article}
\usepackage{amsmath}
\usepackage{amsthm}
\usepackage{microtype}
\newtheorem{theorem}{Theorem}
\begin{document}
\section{Lists}
\begin{itemize}
\item Text after an item, long enough to wrap onto a second line of the list so that
both margins are reached by the item text.
\item \textbf{Bold} start, and \emph{then} more.

A second paragraph in the same item, starting with a ``quote'' and running on.
\item {Grouped} start.
\end{itemize}
\begin{enumerate}
\item Visible text. \item[(x)] Label given.
\end{enumerate}
\begin{description}
\item[Term] Wonderful description text.
\end{description}
\section{Theorems and quotes}
\begin{theorem}
Your theorem text, which also wraps over the end of the first line to show the
protrusion at the margins of a theorem.
\end{theorem}
\begin{quote}
Quoted text, with a final full stop.
\end{quote}
\begin{center}
Centred text.
\end{center}
Text with a footnote.\footnote{The footnote text starts here.}
\section{Equations}
\begin{equation}
a = b
\end{equation}
\begin{align}
c &= d \\
e &= f \tag{T}
\end{align}
See \eqref{x} and the verbatim:
\begin{verbatim}
verbatim text, not protruded.
\end{verbatim}
\begin{equation}\label{x}
g = h
\end{equation}
\end{document}
